\documentclass[11pt,a4paper]{article}
\usepackage[T1]{fontenc}
\usepackage[utf8]{inputenc}
\usepackage[main=italian,provide=*]{babel}
\usepackage{amsmath,amssymb}
\usepackage{geometry}
\geometry{margin=2.3cm,top=2.5cm,bottom=2.7cm}
\usepackage{booktabs}
\usepackage[table]{xcolor}
\usepackage{tcolorbox}
\tcbuselibrary{skins,breakable}
\usepackage{titlesec}
\usepackage{enumitem}
\usepackage{seqsplit}
\usepackage{needspace}
\usepackage{fancyhdr}
\usepackage{hyperref}
\hypersetup{colorlinks=true,linkcolor=Sepia,citecolor=Sepia,urlcolor=Sepia}

% --- palette ---
\definecolor{inkblue}{HTML}{1B3A5C}
\definecolor{lightblue}{HTML}{EAF1F8}
\definecolor{accent}{HTML}{B0562A}
\definecolor{softgray}{HTML}{6B6B6B}
\definecolor{okgreen}{HTML}{2E6B3E}
\definecolor{okgreenbg}{HTML}{EAF5EC}
\definecolor{notebar}{HTML}{9C6B12}

% --- titoli di sezione ---
\titleformat{\section}
  {\normalfont\Large\bfseries\color{inkblue}}
  {\thesection}{0.6em}{}[{\color{inkblue!40}\titlerule}]
\titlespacing{\section}{0pt}{0.8em}{0.4em}
\titleformat{\subsection}
  {\normalfont\large\bfseries\color{accent}}
  {}{0em}{}
\titlespacing{\subsection}{0pt}{0.5em}{0.2em}

\pagestyle{fancy}
\fancyhf{}
\renewcommand{\headrulewidth}{0.4pt}
\fancyhead[L]{\small\color{softgray}RBS vs $W_{ij}(\theta)$: scelta del blocco $M$-conservante}
\fancyhead[R]{\small\color{softgray}\thepage}
\fancyfoot[C]{\small\color{softgray}Tesi triennale, Samuele Schiavi --- Relatore: Prof.\ Alessandro Chiesa}

% --- box riutilizzabili ---
\newtcolorbox{keyeq}[1][]{
  colback=lightblue, colframe=inkblue!70, boxrule=0.6pt,
  arc=2pt, left=8pt, right=8pt, top=4pt, bottom=4pt, #1
}
\newtcolorbox{panelbox}[1]{
  colback=white, colframe=softgray!50, boxrule=0.5pt,
  arc=2pt, left=7pt, right=7pt, top=3pt, bottom=3pt,
  fonttitle=\bfseries\color{inkblue}, title={#1}
}
\newtcolorbox{okbox}[1]{
  colback=okgreenbg, colframe=okgreen!60, boxrule=0.6pt,
  arc=2pt, left=7pt, right=7pt, top=3pt, bottom=3pt,
  fonttitle=\bfseries\color{okgreen}, title={#1}
}
\newtcolorbox{editnote}{
  colback=white, colframe=white, boxrule=0pt,
  borderline west={2.2pt}{0pt}{notebar},
  arc=0pt, left=10pt, right=4pt, top=2pt, bottom=2pt,
  fontupper=\itshape\small\color{softgray!90!black}
}

\title{\vspace{-1.5em}\color{inkblue}RBS vs $W_{ij}(\theta)$\\[0.1em]
\Large\normalfont\color{softgray}Perché il blocco $M$-conservante del PMA-2q resta RBS anche per il trimero}
\author{Note di lavoro --- Tesi triennale, Samuele Schiavi\\ Relatore: Prof.\ Alessandro Chiesa}
\date{}

\begin{document}
\sloppy
\maketitle
\thispagestyle{fancy}
\vspace{-2em}

\textbf{Contesto.} In \texttt{confronto\_ansatz\_entangler.ipynb} ($N=2$) il
blocco $M$-conservante del PMA è realizzato come rotazione di Givens
reale (RBS). Crippa et al.\ (\emph{Magnetochemistry} 7, 117, 2021), nel
costruire il PMA per anelli di spin a $N=4,6$, usano invece un gate
diverso, $W_{ij}(\theta)=e^{-i\theta\,\mathbf s_i\cdot\mathbf s_j}$.
Questo documento deriva esplicitamente la differenza fra i due gate,
quantifica il costo in gate e l'espressività, e motiva la scelta di
\textbf{restare su RBS} anche per l'estensione a $N=3$ (triangolo
isoscele).

\section{Il gate RBS}

Nel sottospazio a due qubit, ristretto al settore $M=0$ (cioè allo
spazio $\{|01\rangle,|10\rangle\}$), RBS è la rotazione di Givens reale:
\begin{keyeq}
\centering
$\mathrm{RBS}(\varphi)=\begin{pmatrix}1&0&0&0\\0&\cos\varphi&-\sin\varphi&0\\0&\sin\varphi&\cos\varphi&0\\0&0&0&1\end{pmatrix}$
\quad (base $|00\rangle,|01\rangle,|10\rangle,|11\rangle$)
\end{keyeq}
Identità su $|00\rangle,|11\rangle$; sul sottospazio
$\{|01\rangle,|10\rangle\}$ agisce come $R_y(2\varphi)$ reale (a meno di
parametrizzazione dell'angolo). Applicato a $|01\rangle$:
$$\mathrm{RBS}(\varphi)\,|01\rangle=\cos\varphi\,|01\rangle+\sin\varphi\,|10\rangle.$$
\textbf{Punto chiave:} i coefficienti sono \textbf{reali} per ogni
$\varphi$.

Implementazione già validata nel progetto (self-test $\sim10^{-16}$
contro il target):
$$\mathrm{RBS}(\varphi)=(H\otimes H)\;\mathrm{CZ}\;\big(R_y(\varphi)_{q_0}\otimes R_y(-\varphi)_{q_1}\big)\;\mathrm{CZ}\;(H\otimes H).$$
Costo: \textbf{2 gate a due qubit} ($CZ$, equivalenti a CNOT tramite
$CZ=(I\otimes H)\,\mathrm{CNOT}\,(I\otimes H)$).

\section{Il gate $W_{ij}(\theta)$ (Crippa et al.)}

$$W_{ij}(\theta)=e^{-i\theta\,\mathbf s_i\cdot\mathbf s_j}, \qquad \mathbf s=\boldsymbol\sigma/2.$$

\subsection*{Spettro di $\mathbf s_i\cdot\mathbf s_j$}

Per due spin-1/2, usando
$\mathbf s_i\cdot\mathbf s_j=\tfrac12(\mathbf S^2-s_i^2-s_j^2)$ con
$s_i^2=s_j^2=\tfrac34$:
\begin{itemize}[leftmargin=1.3em,itemsep=0.1em]
\item \textbf{Tripletto} ($S=1$): $\mathbf s_i\cdot\mathbf s_j=\tfrac12(2-\tfrac34-\tfrac34)=\tfrac14$;
\item \textbf{Singoletto} ($S=0$): $\mathbf s_i\cdot\mathbf s_j=\tfrac12(0-\tfrac34-\tfrac34)=-\tfrac34$.
\end{itemize}
Quindi $W_{ij}(\theta)$ è \textbf{diagonale nella base
singoletto/tripletto}, con autovalori $e^{-i\theta/4}$ (tripletto) ed
$e^{i3\theta/4}$ (singoletto) --- una fase relativa fra i due settori,
non una rotazione reale.

\subsection*{Azione esplicita su $|01\rangle$}

Scrivendo $|01\rangle$ nella base singoletto/tripletto e applicando
$W_{ij}(\theta)$, dopo aver raccolto la fase globale $e^{-i\theta/4}$:
$$W_{ij}(\theta)\,|01\rangle=\frac{e^{-i\theta/4}}{2}\Big[(1+e^{i\theta})|01\rangle+(1-e^{i\theta})|10\rangle\Big].$$

\subsection*{Confronto diretto con RBS: stesse probabilità, fase diversa}

Usando $1+e^{i\theta}=2\cos(\theta/2)e^{i\theta/2}$ e
$1-e^{i\theta}=-2i\sin(\theta/2)e^{i\theta/2}$, a meno della fase
globale $e^{i\theta/4}$ e ponendo $\varphi=\theta/2$:

\begin{okbox}{Risultato del confronto}
$$W_{ij}(2\varphi)\,|01\rangle \;\dot=\; \cos\varphi\,|01\rangle-i\sin\varphi\,|10\rangle
\qquad\text{contro}\qquad
\mathrm{RBS}(\varphi)\,|01\rangle=\cos\varphi\,|01\rangle+\sin\varphi\,|10\rangle.$$
\end{okbox}

Le \textbf{probabilità} di misurare $|01\rangle$ o $|10\rangle$ sono
identiche in entrambi i casi ($\cos^2\varphi$, $\sin^2\varphi$): i due
gate sono equivalenti \emph{a livello di popolazioni}. Quello che li
distingue è la \textbf{fase relativa} fra le due ampiezze: reale ($+1$,
riassorbibile nel segno) per RBS, \textbf{immaginaria} ($-i$) per $W$.
Non è un dettaglio di convenzione: è una proprietà strutturale del
generatore. $\mathbf s_i\cdot\mathbf s_j$, ristretto al sottospazio
$\{|01\rangle,|10\rangle\}$, è proporzionale a una matrice di Pauli $X$
effettiva (scambio puro), il cui esponenziale genera rotazioni con
off-diagonal immaginario --- la stessa famiglia del gate iSWAP, non
quella del gate SWAP/Givens reale. RBS, per costruzione ($R_y$ reali
coniugati da Hadamard/CZ), resta nella famiglia delle rotazioni
\textbf{reali}.

\needspace{10\baselineskip}
\section{Costo in gate}

\begin{center}
\renewcommand{\arraystretch}{1.3}
\begin{tabular}{@{}lp{5.2cm}c@{}}
\toprule
\rowcolor{inkblue!10}
\textbf{Gate} & \textbf{Decomposizione} & \textbf{Gate a 2 qubit}\\
\midrule
$\mathrm{RBS}(\varphi)$ & $H\text{-}CZ\text{-}R_y\text{-}R_y\text{-}CZ\text{-}H$ & \textbf{2} ($CZ$)\\
$W_{ij}(\theta)$ (Crippa, eq.\ 6) & sequenza a gate singolo qubit + CNOT & \textbf{3} (CNOT)\\
\bottomrule
\end{tabular}
\end{center}

Sostituire RBS con $W$ \textbf{non} rende l'ansatz più compatto: lo
appesantisce di un gate a due qubit per blocco (50\% in più).

\begin{editnote}
Nota per evitare un'ambiguità di lettura. Questo costo aggiuntivo è un
argomento \emph{a favore} di RBS, non a favore di $W$ --- vale ancora
di più, non di meno, quando in Parte 2 si introduce il rumore: ogni
gate a due qubit (CNOT o CZ) è un canale di errore in più
(depolarizzazione, errori di gate, ecc.\ --- coerente con l'analisi già
fatta in \texttt{costo\_gate\_trotter.md} per il blocco Trotter del
dimero, dove il conteggio di CNOT è stato esplicitamente il criterio
per confrontare le decomposizioni). Meno CNOT per blocco significa un
ansatz più robusto quando si passa al simulatore rumoroso. RBS resta
quindi la scelta giusta anche in Parte 2, a maggior ragione, non solo
in Parte 1.
\end{editnote}

Va sottolineato che il passaggio a Parte 2 cambia la \emph{dinamica}
(stato puro $\to$ operatore densità, evoluzione unitaria $\to$ canali
di Kraus/Lindblad), non l'\emph{Hamiltoniana}: $H$ resta la stessa
matrice reale simmetrica. L'argomento di realtà della sez.~4 (teorema
spettrale $\Rightarrow$ ground state reale $\Rightarrow$ RBS
espressivamente sufficiente) dipende solo da $H$, non da come si
esegue la ricerca variazionale né da quanto rumore è presente nel
circuito. Il rumore (bit-flip, $T_1$, $T_2$, ecc.) non richiede fasi
complesse nell'\emph{ansatz} per essere modellato: agisce come canale
separato sull'output del circuito, non sui parametri variazionali. Le
due motivazioni --- realtà di $H$ e minor costo in gate --- si
sommano, non si sostituiscono a vicenda.

\section{Espressività: perché la fase immaginaria è superflua qui}

\textbf{Argomento di realtà (teorema spettrale).} L'Hamiltoniana del
trimero isoscele,
$$H=J\,\boldsymbol\sigma_1\cdot\boldsymbol\sigma_2+J'(\boldsymbol\sigma_2\cdot\boldsymbol\sigma_3+\boldsymbol\sigma_3\cdot\boldsymbol\sigma_1)+b\sum_iZ_i,$$
è \textbf{reale simmetrica} in ogni base computazionale:
$\boldsymbol\sigma_i\cdot\boldsymbol\sigma_j=X_iX_j+Y_iY_j+Z_iZ_j$ ha
matrice reale (anche $Y_iY_j$: il prodotto di due entrate immaginarie
di $Y$ dà un'entrata reale --- stesso argomento già verificato per il
dimero), e $\sum_iZ_i$ è diagonale reale. Per il teorema spettrale, una
matrice reale simmetrica ammette una base ortonormale di
\textbf{autovettori reali}: lo stato fondamentale può sempre essere
scelto reale, qualunque sia la degenerazione.

\textbf{Conseguenza per l'ansatz.} Un ansatz che deve raggiungere solo
stati reali non ha bisogno di esplorare fasi complesse. $W_{ij}(\theta)$
esplora l'intero $U(2)$ ristretto al settore $M=0$ della coppia
(include la fase $-i$ di iSWAP); RBS esplora solo il sottogruppo reale
$SO(2)$. La differenza è esattamente lo stesso tipo di ridondanza già
isolata e quantificata per gli ansatz A--D nel dimero (dove
$n_\text{par}-3$ parametri risultavano gauge, non fisici): qui il
"parametro extra" non è un angolo in più ma un \textbf{grado di libertà
di fase per blocco}, sempre sprecato.

\textbf{Quando smetterebbe di essere superfluo.} Il vincolo di realtà
cade se si introduce un termine che rende $H$ genuinamente complessa.
Il DM già usato nel dimero, nella forma $D(X_iZ_j-Z_iX_j)$, resta
\textbf{reale} (stessa argomentazione) --- quindi anche estendendolo al
triangolo con la stessa forma, RBS continuerebbe a bastare. Servirebbe
passare a $W$ (o equivalentemente aggiungere $R_z$ nel blocco) solo con
un termine genuinamente immaginario, come il DM "canonico"
$D(X_iY_j-Y_iX_j)$ --- esplicitamente fuori dallo scope attuale (vedi
punto 9 di \texttt{domande\_relatore.md}).

\needspace{12\baselineskip}
\section{Verifica numerica di quanto derivato sopra}

Controllo diretto (non solo simbolico) che $W_{ij}(2\varphi)$ e
$\mathrm{RBS}(\varphi)$ diano le stesse probabilità ma fasi diverse,
per alcuni valori di $\varphi$:

\begin{panelbox}{Codice di verifica}
\small
\begin{verbatim}
import numpy as np

def sdot(theta):
    """W_ij(theta) = exp(-i*theta*s_i.s_j), spin s=sigma/2,
       base |00>,|01>,|10>,|11>."""
    X = np.array([[0,1],[1,0]], dtype=complex)
    Y = np.array([[0,-1j],[1j,0]], dtype=complex)
    Z = np.array([[1,0],[0,-1]], dtype=complex)
    s = [0.5*X, 0.5*Y, 0.5*Z]
    sdotop = sum(np.kron(a,a) for a in s)
    w, v = np.linalg.eigh(sdotop)
    return v @ np.diag(np.exp(-1j*theta*w)) @ v.conj().T

def rbs(phi):
    c, s = np.cos(phi), np.sin(phi)
    return np.array([[1,0,0,0],[0,c,-s,0],[0,s,c,0],[0,0,0,1]],
                     dtype=complex)

ket01 = np.array([0,1,0,0], dtype=complex)
for phi in [0.3, 0.9, 1.4]:
    out_W   = sdot(2*phi) @ ket01
    out_RBS = rbs(phi) @ ket01
    print(f"phi={phi:.2f}  W: {out_W.round(4)}   "
          f"RBS: {out_RBS.round(4)}  "
          f"|prob uguali|: "
          f"{np.allclose(np.abs(out_W)**2, np.abs(out_RBS)**2)}")
\end{verbatim}
\end{panelbox}

Atteso (coerente con la derivazione della sez.~2.3): componente su
$|01\rangle$ reale e uguale in entrambi i casi ($\cos\varphi$);
componente su $|10\rangle$ reale positiva per RBS ($\sin\varphi$),
immaginaria negativa per $W$ ($-i\sin\varphi$, a meno della fase
globale $e^{i\theta/4}$) --- probabilità identiche, fase
strutturalmente diversa.

\needspace{10\baselineskip}
\section{Conclusione operativa}

\begin{enumerate}[leftmargin=1.4em,itemsep=0.3em]
\item \textbf{RBS resta il blocco $M$-conservante di riferimento},
anche per il trimero: stesso potere espressivo di $W$ sulle
popolazioni, costo inferiore (2 vs 3 gate a due qubit per blocco),
nessuna capacità aggiuntiva sprecata dato che $H$ è reale.
\item Quello che vale la pena \textbf{riusare concettualmente} dal
paper di Crippa et al.\ non è il gate, ma la \textbf{struttura}: layer
su più legami (qui i tre bond $12,23,31$), stato iniziale scelto nel
settore di simmetria del multipletto target (qui, guidato dalla
decomposizione di Kambe: $S_{12}=0$ per il blocco C, $S_{12}=1$ per
A/B), e l'euristica sull'ordine dei legami (applicare prima il bond che
porta lo stato locale singoletto/tripletto).
\item Il blocco RBS applicato su una coppia qualunque, con il terzo
qubit spettatore, \textbf{conserva comunque $M$ totale} a 3 qubit
(conserva $m_i+m_j$ della coppia, lascia $m_k$ invariato): la proprietà
chiave del PMA si eredita intatta nel passaggio a $N=3$.
\item Il vincolo va \textbf{ridiscusso} solo se in futuro si introduce
un termine DM immaginario sul triangolo (non in scope ora).
\end{enumerate}

\section{Perché RBS trova il ground state anche sopra soglia ($b/J>2$)}

\subsection*{Il problema: RBS da solo conserva $M$}

Un singolo blocco RBS agisce solo sul settore $\{|01\rangle,|10\rangle\}$
e lascia invariati $|00\rangle$ e $|11\rangle$. Partendo da $|01\rangle$:
$$\mathrm{RBS}(\varphi)|01\rangle=\cos\varphi\,|01\rangle+\sin\varphi\,|10\rangle,$$
cioè lo stato rimane sempre a $M=0$. Il ground state sopra soglia
($b/J>2$) è $|t_-\rangle=|11\rangle$ con $M=-1$: irraggiungibile
applicando solo RBS.

Questa è la stessa ragione per cui il PMA di base (1 parametro, solo un
blocco RBS dopo lo stato iniziale $X|00\rangle=|10\rangle$) fallisce a
$b/J>2$ anche con $D=0$: non è un effetto del termine DM, è
semplicemente l'assenza di un meccanismo per uscire dal settore $M=0$.

\subsection*{La soluzione: gli strati $R_y$ rompono la conservazione di $M$}

L'ansatz PMA-2q non è solo RBS. Ogni strato ha la struttura
$$\bigl(R_y(\alpha_k)\otimes R_y(\beta_k)\bigr)\cdot\mathrm{RBS}(\varphi_k),$$
con $\alpha_k\neq\beta_k$ in generale. Un gate $R_y(\alpha)$ su un
singolo qubit $q_i$ agisce localmente come
$$R_y(\alpha)|0\rangle=\cos\tfrac\alpha2|0\rangle-\sin\tfrac\alpha2|1\rangle,\quad
R_y(\alpha)|1\rangle=\sin\tfrac\alpha2|0\rangle+\cos\tfrac\alpha2|1\rangle,$$
e quindi \textbf{mescola $|0\rangle$ e $|1\rangle$ su quel qubit}.
Applicato dopo RBS (che ha prodotto uno stato in $M=0$), un $R_y$
locale redistribuisce ampiezza verso settori $M=\pm1$, incluso
$|11\rangle$. La sequenza completa esplora tutto lo spazio reale dei 2
qubit, non solo il settore $M=0$.

\textbf{Meccanismo esplicito.} Dopo il blocco RBS lo stato è
$|\psi\rangle=a|01\rangle+b|10\rangle$ (con $a,b$ reali,
$a^2+b^2=1$). Applicando $R_y(\alpha)\otimes R_y(\beta)$, il
coefficiente di $|11\rangle$ nel risultato è
$$-a\sin\tfrac\alpha2\cos\tfrac\beta2+b\cos\tfrac\alpha2\sin\tfrac\beta2,$$
che è zero solo se $\alpha=\beta$ oppure $a=b=0$. Per $\alpha\neq\beta$
generici lo stato ha componente non nulla su $|11\rangle$ (settore
$M=-1$): l'ottimizzatore può quindi portare $E(\boldsymbol\theta)$
verso l'energia di $|11\rangle$ sopra soglia.

\subsection*{Perché $W$ non aveva questo problema (ma per una ragione diversa da quella apparente)}

$W_{ij}(\theta)=e^{-i\theta\,\mathbf s_i\cdot\mathbf s_j}$
\textbf{conserva anch'esso $M$} (è diagonale nella base
$\{$singoletto, tripletto$\}$, che è una base con $M$ definito). La
ragione per cui il PMA originale di Crippa et al.\ funzionava sopra
soglia non è che $W$ esce dal settore $M=0$, ma che
nell'implementazione di Crippa il gate $W$ era sempre accompagnato da
rotazioni locali, esattamente come nel PMA-2q qui --- quelle rotazioni
locali sono il meccanismo che popola $|11\rangle$, non $W$ in sé.

\subsection*{Tabella riassuntiva}

\begin{center}
\renewcommand{\arraystretch}{1.35}
\begin{tabular}{@{}p{3.6cm}cc p{4.4cm}@{}}
\toprule
\rowcolor{inkblue!10}
\textbf{Componente} & \textbf{Conserva $M$?} & \textbf{Raggiunge $|11\rangle$?} & \textbf{Ruolo}\\
\midrule
$\mathrm{RBS}(\varphi)$ da solo & sì ($M{=}0$) & no & entangler nel settore $M=0$\\
$R_y(\alpha)\otimes R_y(\alpha)$ (par.\ condiviso) & no & sì & non abbastanza per $\mathcal F{=}1$ (vedi \texttt{\seqsplit{analisi\_rotazioni\_PMA.ipynb}})\\
$R_y(\alpha)\otimes R_y(\beta)$ (par.\ indipendenti) & no & sì & popola tutti i settori, necessario per $\mathcal F{=}1$\\
PMA-2q$\cdot$3 ($K{=}1$, 3 par.\ totali) & no & sì & $\mathcal F{=}1$ su tutto il range col numero minimo teorico di parametri\\
\bottomrule
\end{tabular}
\end{center}

\begin{editnote}
Nota sulla ridondanza del parametro condiviso. Se $\alpha=\beta$
(stessa rotazione su entrambi i qubit), il coefficiente di $|11\rangle$
diventa $\sin\tfrac\alpha2\cos\tfrac\alpha2(b-a)=\tfrac12\sin\alpha\,(b-a)$,
che si annulla per $a=b$ (mescolamento massimo, cioè proprio $b/J=2$,
il punto critico). Questo spiega quantitativamente perché un parametro
condiviso non basta: fallisce esattamente nel punto più difficile, dove
il mescolamento singoletto/$|t_-\rangle$ è massimo e il ground state
richiede componenti uguali di entrambi.
\end{editnote}

\bigskip
\noindent\textbf{Riferimento.} Dimostrazione completa con tutti i casi
che falliscono: \texttt{analisi\_rotazioni\_PMA.ipynb}.

\end{document}
